\documentclass[12pt,a4paper]{article}
\usepackage[utf8]{inputenc}
\usepackage[T1]{fontenc}
\usepackage{lmodern}
\usepackage{amsmath,amssymb,amsfonts,amsthm}
\usepackage{mathtools}
\usepackage{geometry}
\geometry{top=1in, bottom=1in, left=1in, right=1in}
\usepackage{xcolor}
\usepackage{booktabs}
\usepackage{fancyhdr}
\usepackage{hyperref}

\definecolor{navy}{RGB}{20, 45, 85}

\hypersetup{
    colorlinks=true,
    linkcolor=navy,
    citecolor=navy,
    urlcolor=navy
}

\theoremstyle{definition}
\newtheorem{task}{Task}[section]
\newtheorem{solutionstep}{Step}[subsection]

\theoremstyle{plain}
\newtheorem{mytheorem}{Theorem}[section]
\newtheorem{mylemma}[mytheorem]{Lemma}

\pagestyle{fancy}
\fancyhf{}
\lhead{\small\sffamily Assignment Submission | MATH 402}
\rhead{\small\sffamily Student: Chandramani Verma}
\lfoot{\small\sffamily Roll No: 2026-MATH-042}
\rfoot{\small\sffamily Page \thepage}
\renewcommand{\headrulewidth}{0.4pt}
\renewcommand{\footrulewidth}{0.4pt}

\setlength{\parskip}{0.6em}
\setlength{\parindent}{0pt}

\begin{document}

% ----------------------------------------------------
% OFFICIAL ASSIGNMENT COVER PAGE
% ----------------------------------------------------
\begin{titlepage}
    \centering
    \vspace*{1.5cm}
    
    {\Large \scshape Department of Mathematics \& Control Engineering}\\[0.3cm]
    {\large M.Sc. Applied Mathematics Program}\\[1.5cm]
    
    \rule{\linewidth}{1.2pt}\\[0.4cm]
    {\huge \bfseries Coursework Assignment 1: State-Space Modeling, Autonomous Dynamics \& Linear Systems}\\[0.3cm]
    \rule{\linewidth}{1.2pt}\\[1.8cm]
    
    \begin{minipage}{0.48\textwidth}
        \begin{flushleft} \large
        \textbf{Student Details:}\\
        Name: Chandramani Verma\\
        Roll Number: \texttt{2026-MATH-042}\\
        Semester: II (2025--2026)
        \end{flushleft}
    \end{minipage}
    ~
    \begin{minipage}{0.48\textwidth}
        \begin{flushright} \large
        \textbf{Course & Submission Info:}\\
        Course Code: \textbf{MATH 402}\\
        Course Title: Dynamical Systems\\
        Date of Submission: June 28, 2026
        \end{flushright}
    \end{minipage}

    \vfill
    
    {\small \scshape Prepared independently for academic course evaluation.}
    \vspace*{1cm}
\end{titlepage}

\newpage
\tableofcontents
\newpage

% ----------------------------------------------------
% TASK 1
% ----------------------------------------------------
\section{Task 1: Formulation of Discrete and Continuous State-Space Models}

In this first task, I formulate mathematical state-space representations for discrete and continuous physical processes.

\subsection{My Understanding of State Vectors}
When I analyze a physical or economic system, I define its \textbf{state vector} $x(t) \in \mathbb{R}^n$ as the minimum set of variables required to fully describe the system's condition at time $t$. If I know $x(t)$ at the present moment, I can determine all future behavior using the state transition law without needing past history.

\subsection{Physical Application: Damped Harmonic Oscillator}
To demonstrate continuous-time modeling, I consider a mechanical particle moving under linear spring restoration and friction. Let $p(t)$ be its position and $v(t) = \dot{p}(t)$ be its velocity. 

Position $p(t)$ alone is insufficient to predict future motion because the particle's direction of movement is unknown. Therefore, I construct the 2D state vector:
\[
x(t) = \begin{bmatrix} p(t) \\ v(t) \end{bmatrix} \in \mathbb{R}^2
\]
The second-order Newton law $\ddot{p} + \beta \dot{p} + k_0 p = 0$ converts into my first-order matrix differential system:
\[
\begin{bmatrix} \dot{p}(t) \\ \dot{v}(t) \end{bmatrix} = \begin{bmatrix} 0 & 1 \\ -k_0 & -\beta \end{bmatrix} \begin{bmatrix} p(t) \\ v(t) \end{bmatrix}
\]
Here, $A = \begin{bmatrix} 0 & 1 \\ -k_0 & -\beta \end{bmatrix}$ represents my system matrix.

\subsection{Discrete-Time Financial Recurrence}
For discrete-time dynamics, time progresses in integer steps $k \in \{0, 1, 2, \dots\}$. The evolution obeys a difference equation $x(k+1) = f(x(k))$.

\textbf{My Numerical Example:} I evaluate an investment account started with an initial capital of $x(0) = \$3,200$ earning an annual interest rate of $4.8\%$, compounded yearly. 
\[
x(k+1) = 1.048 \, x(k)
\]
Iterating from my initial value $x_0 = 3200$:
\begin{align*}
x(1) &= 1.048 \times 3200 = 3353.60 \\
x(2) &= 1.048 \times 3353.60 = (1.048)^2 \times 3200 = 3514.57 \\
x(k) &= (1.048)^k \cdot 3200
\end{align*}
This closed formula allows me to compute the exact account balance at any year $k$.

\newpage
% ----------------------------------------------------
% TASK 2
% ----------------------------------------------------
\section{Task 2: Comparative Analysis of Autonomous and Non-Autonomous Dynamics}

In this task, I compare continuous dynamical systems based on explicit time dependence.

\subsection{Autonomous Systems ($\dot{x} = f(x)$)}
I classify a continuous system as \textbf{autonomous} when the vector field $f(x)$ depends solely on the state vector $x(t)$, with no explicit appearance of time $t$:
\[
\frac{dx}{dt} = f(x)
\]

\textbf{Key Behavior I Observed:}
\begin{enumerate}
    \item \textbf{Time Invariance:} The vector field stays constant in time. Shifting my starting time from $t=0$ to $t=t_0$ shifts the solution curve in time without altering its trajectory shape.
    \item \textbf{Fixed Equilibria:} I determine equilibrium points $x^*$ by solving $f(x^*) = 0$. At these points, $\dot{x} = 0$, so the system stays stationary.
    \item \textbf{No Trajectory Crossings:} In a 2D phase portrait, solution curves never intersect one another due to existence-uniqueness properties.
\end{enumerate}

\subsection{Non-Autonomous Systems ($\dot{x} = f(x, t)$)}
Conversely, a system is \textbf{non-autonomous} if time $t$ appears explicitly in the derivative function:
\[
\frac{dx}{dt} = f(x, t)
\]

\textbf{Key Behavior I Observed:}
\begin{enumerate}
    \item \textbf{Time-Varying Vector Fields:} The velocity field changes continuously over time, as seen in forced pendulums or seasonal biological models.
    \item \textbf{Intersecting Spatial Projections:} 2D spatial projections of trajectories can cross because the vector field changes direction at identical points at different times.
\end{enumerate}

\begin{table}[h]
\centering
\caption{My Summary Table: Autonomous vs. Non-Autonomous Systems}
\vspace{0.3em}
\begin{tabular}{lll}
\toprule
\textbf{Feature} & \textbf{Autonomous ($\dot{x} = f(x)$)} & \textbf{Non-Autonomous ($\dot{x} = f(x,t)$)} \\
\midrule
Explicit Time Variable & Absent & Present \\
Vector Field & Constant across time & Time-varying \\
Equilibrium Points & Fixed roots of $f(x^*)=0$ & Moving roots of $f(x^*,t)=0$ \\
Phase Trajectories & Non-crossing in 2D & Projections may intersect \\
\bottomrule
\end{tabular}
\end{table}

\newpage
% ----------------------------------------------------
% TASK 3
% ----------------------------------------------------
\section{Task 3: Linear State-Space Analysis and Exact Matrix Exponential Solutions}

In this task, I investigate linear homogeneous systems $\dot{x}(t) = A x(t)$ with initial condition $x(0) = x_0$.

\subsection{Uncoupled 2D Linear System}
When $A$ is a diagonal matrix $A = \text{diag}(\lambda_1, \lambda_2)$, the system uncouples into independent scalar ODEs.

\textbf{My Uncoupled Example:} I consider the planar system:
\[
\begin{cases}
\dot{x}_1 = -2 x_1 \\
\dot{x}_2 = 3 x_2
\end{cases}
\implies A = \begin{bmatrix} -2 & 0 \\ 0 & 3 \end{bmatrix}
\]
Solving each equation independently gives:
\[
x_1(t) = c_1 e^{-2t}, \qquad x_2(t) = c_2 e^{3t}
\]
To find the phase plane curves, I express $e^t$ in terms of $x_1$: $e^{-t} = (x_1 / c_1)^{1/2} \implies e^t = (c_1 / x_1)^{1/2}$. Substituting into $x_2$:
\[
x_2(t) = c_2 \left( \frac{c_1}{x_1(t)} \right)^{3/2} = \frac{C_0}{x_1^{1.5}}, \quad \text{where } C_0 = c_2 c_1^{1.5}
\]
As $t \to \infty$, $x_1(t) \to 0$ and $x_2(t) \to \pm \infty$. The origin $(0,0)$ is an \textbf{unstable saddle point}.

\subsection{Matrix Exponential and Solution Uniqueness}
For general square matrices $A \in \mathbb{R}^{n \times n}$, I define the matrix exponential via the power series:
\[
e^{At} = \sum_{j=0}^{\infty} \frac{A^j t^j}{j!} = I + A t + \frac{A^2 t^2}{2!} + \frac{A^3 t^3}{3!} + \cdots
\]

\textbf{My Uniqueness Proof Check:} To prove that $x(t) = e^{At} x_0$ is the unique solution to $\dot{x} = A x, x(0) = x_0$:
\begin{enumerate}
    \item I differentiate $e^{At}$: $\frac{d}{dt} e^{At} = A e^{At}$. Thus $x'(t) = A e^{At} x_0 = A x(t)$, so it satisfies the differential equation.
    \item Let $z(t)$ be any other solution. I define $w(t) = e^{-At} z(t)$.
    \item Differentiating $w(t)$: $w'(t) = -A e^{-At} z(t) + e^{-At} z'(t) = -A e^{-At} z(t) + e^{-At} A z(t) = 0$.
    \item Since $w'(t) = 0$, $w(t)$ is constant. At $t=0$, $w(0) = z(0) = x_0$. Therefore $w(t) = x_0 \implies z(t) = e^{At} x_0$. This confirms uniqueness!
\end{enumerate}

\subsection{My Fully Worked Numerical Problem}
\textbf{Problem Statement:} Solve the coupled linear system:
\[
\begin{cases}
\dot{x}(t) = 3 x(t) + 2 y(t) \\
\dot{y}(t) = 4 x(t) + 1 y(t)
\end{cases}
\quad \text{with initial condition } x(0) = 4, \, y(0) = -2
\]

\textbf{Step 1: Characteristic Equation and Eigenvalues}
Writing in matrix form $\dot{X} = A X$ with $A = \begin{bmatrix} 3 & 2 \\ 4 & 1 \end{bmatrix}$:
\[
\det(A - \lambda I) = \det \begin{bmatrix} 3-\lambda & 2 \\ 4 & 1-\lambda \end{bmatrix} = (3-\lambda)(1-\lambda) - 8 = \lambda^2 - 4\lambda - 5 = 0
\]
Factoring gives $(\lambda - 5)(\lambda + 1) = 0$, yielding eigenvalues:
\[
\lambda_1 = 5, \qquad \lambda_2 = -1
\]

\textbf{Step 2: Eigenvector Calculation}
\begin{itemize}
    \item For $\lambda_1 = 5$: $(A - 5I) v_1 = 0 \implies \begin{bmatrix} -2 & 2 \\ 4 & -4 \end{bmatrix} \begin{bmatrix} v_{11} \\ v_{12} \end{bmatrix} = 0 \implies v_1 = \begin{bmatrix} 1 \\ 1 \end{bmatrix}$.
    \item For $\lambda_2 = -1$: $(A + I) v_2 = 0 \implies \begin{bmatrix} 4 & 2 \\ 4 & 2 \end{bmatrix} \begin{bmatrix} v_{21} \\ v_{22} \end{bmatrix} = 0 \implies v_2 = \begin{bmatrix} 1 \\ -2 \end{bmatrix}$.
\end{itemize}

\textbf{Step 3: Modal Matrix and Inverse}
\[
P = \begin{bmatrix} 1 & 1 \\ 1 & -2 \end{bmatrix}, \quad \det(P) = -3 \implies P^{-1} = \frac{1}{3} \begin{bmatrix} 2 & 1 \\ 1 & -1 \end{bmatrix}
\]

\textbf{Step 4: State Solution $X(t) = P e^{Dt} P^{-1} X_0$}
First, compute transformed initial state $C_0 = P^{-1} X_0$:
\[
C_0 = \frac{1}{3} \begin{bmatrix} 2 & 1 \\ 1 & -1 \end{bmatrix} \begin{bmatrix} 4 \\ -2 \end{bmatrix} = \frac{1}{3} \begin{bmatrix} 8 - 2 \\ 4 + 2 \end{bmatrix} = \begin{bmatrix} 2 \\ 2 \end{bmatrix}
\]
Now multiply $X(t) = P \begin{bmatrix} 2 e^{5t} \\ 2 e^{-t} \end{bmatrix}$:
\[
X(t) = \begin{bmatrix} 1 & 1 \\ 1 & -2 \end{bmatrix} \begin{bmatrix} 2 e^{5t} \\ 2 e^{-t} \end{bmatrix} = \begin{bmatrix} 2 e^{5t} + 2 e^{-t} \\[0.4em] 2 e^{5t} - 4 e^{-t} \end{bmatrix}
\]

\textbf{My Final Exact Solution:}
\[
x(t) = 2 e^{5t} + 2 e^{-t}, \qquad y(t) = 2 e^{5t} - 4 e^{-t}
\]
\textit{Initial Condition Verification:} At $t=0$, $x(0) = 2+2 = 4$ and $y(0) = 2-4 = -2$. This matches my initial conditions perfectly!

\newpage
% ----------------------------------------------------
% TASK 4
% ----------------------------------------------------
\section{Task 4: Annotated Course Reading & Literature Review}

In this task, I review key academic literature consulted during this coursework:

\begin{enumerate}
    \item \textbf{Hirsch, M. W., Smale, S., \& Devaney, R. L. (2012).} \textit{Differential Equations, Dynamical Systems, and an Introduction to Chaos} (3rd ed.). Academic Press.
    \begin{quote}
    \small \textit{My Note:} I consulted Chapters 1--3 for rigorous proofs of matrix exponentials and phase plane classifications.
    \end{quote}
    
    \item \textbf{Perko, L. (2001).} \textit{Differential Equations and Dynamical Systems} (3rd ed.). Springer-Verlag.
    \begin{quote}
    \small \textit{My Note:} Used for state-space coordinate transformations and canonical Jordan forms.
    \end{quote}
    
    \item \textbf{Khalil, H. K. (2002).} \textit{Nonlinear Systems} (3rd ed.). Prentice Hall.
    \begin{quote}
    \small \textit{My Note:} Reference for non-autonomous system stability and time-varying vector field properties.
    \end{quote}
    
    \item \textbf{Strogatz, S. H. (2014).} \textit{Nonlinear Dynamics and Chaos}. Westview Press.
    \begin{quote}
    \small \textit{My Note:} Provides valuable physical intuition for phase portraits and mechanical oscillator state vectors.
    \end{quote}
\end{enumerate}

\end{document}
